% =====================================================================
% CHAPTER 8: Results and Discussion (9-chapter RUET MTE template)
% Template source: thesis-or-project-template_1745300785.pdf, Chapter 8.
% Assembled per thesis/RESTRUCTURE_PLAN.md from old mainmatter/ch4.tex:
%   - 8.1 Introduction: NEW; compresses the old "Experimental Setup
%     Summary" into pointer sentences (setup body moved to Chapter 7).
%   - 8.2 Investigation and Analysis of Data: old "Training Dynamics" +
%     "Ablation Studies" (three subsections), tables/figures unchanged.
%   - 8.3 Results: the six result sections as subsections, unchanged.
%   - 8.4 Discussion: old "Discussion" unchanged (incl fig:pareto_ours).
%   - 8.5 Summary: NEW; restates headline numbers already present above.
% EXCLUDED: "Societal, Environmental and Sustainability Aspects"
%   (moves to Chapter 4) and the "Experimental Setup Summary" body
%   (moves to Chapter 7). All \label names unchanged, each used once.
% =====================================================================
\chapter{Results and Discussion}

\section{Introduction}
This chapter reports accuracy, efficiency, and cross-domain behaviour
for the trained model: in-domain on FlyingThings3D against published
methods, zero-shot on Middlebury 2014 against matched-protocol
checkpoints, and on-device latency against the real-time target, plus
the training dynamics and ablation studies behind the design in
Chapter 3. Unless noted, results come from one production run trained
on the full SceneFlow finalpass dataset under the protocol and
configuration of Chapter 7 (Table \ref{tab:train_config}). Accuracy
follows \eqref{eq:metrics}; latency is measured at batch one, 384 by
640; cross-domain behaviour uses Middlebury 2014, which the model
never saw during training.

\FloatBarrier
\section{Investigation and Analysis of Data}
This section covers the training trajectory of the production run and
the controlled studies behind the architecture in Chapter 3, built
from a sequence of studies rather than one design pass. After the
training dynamics, the studies are grouped in three: early exploration
that selected the encoder, refinement scheme, and loss on the earlier
tile chassis; the ablation of the recurrent family that produced the
final model; and training-recipe and efficiency studies on the final
model. All studies follow the matched harness protocol of Chapter 3.

\FloatBarrier
\subsection{Training Dynamics}
Figure \ref{fig:training_curves} shows the training trajectory. Loss
falls sharply during the OneCycle warm-up, then descends steadily as the
learning rate anneals toward its floor. Validation end-point error tracks
the loss closely, bottoming out at 0.79 pixels near step 53,000 and
holding flat between 0.80 and 0.81 pixels for several thousand steps as
the rate nears zero. The best checkpoint is taken at step 53,000, not the
final step 60,000, since further training gives no improvement. Figure
\ref{fig:convergence} shows the same pattern: broad surfaces settle early
while object boundaries and thin structures keep sharpening. The run
survived several cloud preemptions: scheduler and optimizer state are
checkpointed and restored exactly, so a preemption costs at most the
steps since the last validation. The one-cycle schedule's high peak rate
acts as a regularizer mid-training, visible as the sustained loss
variance in Figure \ref{fig:training_curves}; its long cosine descent
then turns that exploration into steady gains, more than half the final
improvement coming after the rate drops below a quarter of its peak. The
per-batch loss wobble in the high-rate phase is the schedule working as
designed, not instability, and the flat validation tail confirms 60,000
steps was about right for this model and dataset.

\begin{figure}[!htbp]
\centering
\includegraphics[width=\textwidth]{fig_4_1_training_curves}
\caption{\textbf{Training dynamics of the 60k-step run.} (a) Training loss with the OneCycle learning-rate schedule. (b) Held-out validation EPE and bad-1 on the fixed 400-pair FlyingThings3D subset at the native axis; the star marks the best checkpoint at step 53k.}
\label{fig:training_curves}
\end{figure}

\begin{figure}[!htbp]
\centering
\includegraphics[width=\textwidth]{fig_4_2_convergence}
\caption{\textbf{Qualitative convergence over training.} Predictions for two held-out FlyingThings3D scenes at five checkpoints (1k to 53k steps) against ground truth; turbo colormap, per-scene normalization. Object boundaries and thin structures sharpen progressively as training advances.}
\label{fig:convergence}
\end{figure}

\FloatBarrier
\subsection{Design-Space Exploration of the Tile Chassis}
An earlier campaign on a smaller tile chassis settled three questions.
Table \ref{tab:ablation_encoder} compared a GhostConv stack against two
truncated detection backbones under a matched 20-pair harness; YOLO26s
won on every metric, 0.528 pixel end-point error against 0.625 for
GhostConv and 0.712 for the smaller YOLO26n, at a two-millisecond latency
cost, and was selected for the final model. Table
\ref{tab:ablation_refinement} tested the iterated coarse-to-fine scheme
against two alternatives: extending refinement to finer scales gained
almost nothing (0.573 against 0.578 pixels) at 37 percent more latency,
while HITNet-style single-pass-per-scale propagation collapsed accuracy
to 0.844 pixels. Iterated refinement at coarse scales is thus the
load-bearing choice, kept in the final model. Table
\ref{tab:ablation_loss} compared nine formulations under the same
harness; the threshold-stack objective with the D1-relative hinge won on
end-point error, bad-2, and D1, and is the ancestor of the production
objective of \eqref{eq:loss}. The clearest negative result:
sequence-weighted multi-scale L1, which imitates the recurrent
literature's iteration-weighted loss without exposing true per-iteration
outputs, was the worst of all nine at 1.005 pixels. Faking a sequence
loss across scales does not work; the final model instead exposes true
multi-scale outputs to the objective.

\begin{table}[!htbp]
\centering
\caption{\textbf{Encoder ablation on the tile chassis.} Matched 20-pair harness, identical steps and seed; latency at $384\times 640$, batch 1, RTX 3050.}
\label{tab:ablation_encoder}
\small
\begin{tabular}{lccccc}
\toprule
Encoder & Params (M) & EPE & bad-1 (\%) & D1-all (\%) & Latency (ms)\\
\midrule
GhostConv stack & 0.538 & 0.625 & 13.31 & 3.54 & 23.5\\
YOLO26n (truncated) & 0.808 & 0.712 & 18.07 & 2.90 & 24.1\\
\textbf{YOLO26s (truncated)} & 2.061 & \textbf{0.528} & \textbf{10.88} & \textbf{2.50} & 25.8\\
\bottomrule
\end{tabular}
\end{table}

\begin{table}[!htbp]
\centering
\caption{\textbf{Refinement-scheme ablation on the tile chassis.} Matched 20-pair harness; the iterated coarse-to-fine scheme against extending refinement to finer scales and against single-pass-per-scale propagation.}
\label{tab:ablation_refinement}
\small
\begin{tabular}{lccccc}
\toprule
Scheme & Params (M) & EPE & bad-1 (\%) & D1-all (\%) & Latency (ms)\\
\midrule
\textbf{iterated coarse-to-fine} & 0.538 & \textbf{0.578} & \textbf{10.72} & 3.44 & \textbf{23.8}\\
+ refinement at 1/2 & 0.576 & 0.573 & 11.12 & \textbf{3.25} & 32.5\\
single pass per scale & 0.487 & 0.844 & 19.30 & 4.79 & 25.2\\
\bottomrule
\end{tabular}
\end{table}

\begin{table}[!htbp]
\centering
\caption{\textbf{Loss-formulation ablation on the tile chassis.} Nine objectives under the matched 20-pair harness; the threshold-stack objective with the D1 hinge wins on EPE, bad-2, and D1 together.}
\label{tab:ablation_loss}
\small
\begin{tabular}{lcccc}
\toprule
Objective & EPE & bad-0.5 (\%) & bad-2 (\%) & D1-all (\%)\\
\midrule
plain L1 & 0.663 & 27.6 & 6.09 & 3.99\\
sequence-weighted multi-scale L1 & 1.005 & 53.5 & 13.25 & 4.04\\
L1 + gradient & 0.618 & \textbf{23.9} & 5.93 & 3.92\\
L1 + bad-1 hinge & 0.763 & 45.8 & 6.21 & 3.65\\
cocktail (L1 + grad + bad-1) & 0.692 & 40.0 & 6.12 & 3.68\\
cocktail, 0.5-px hinge & 0.676 & 38.3 & 6.21 & 3.50\\
threshold stack & 0.672 & 38.0 & 5.97 & 3.45\\
\textbf{threshold stack + D1 hinge} & \textbf{0.591} & 31.5 & \textbf{5.64} & \textbf{3.35}\\
Charbonnier & 0.631 & 31.0 & 5.78 & 3.58\\
\bottomrule
\end{tabular}
\end{table}

\FloatBarrier
\subsection{Architecture Ablation of the Recurrent Family}
The second campaign built the recurrent context family that became the
final model, changing one mechanism at a time on a matched 20-pair
harness (Table \ref{tab:ablation_gev4line}). The base design reached
0.424 pixel end-point error; a learned update gate plus the YOLO26s
encoder brought it to 0.319. Rejected additions: a selective recurrent
unit (0.474), a fresh quarter-resolution initialization (0.455), and
the sharp-tail and high-resolution-refine variants, which gave minor
per-metric gains at added latency without improving the headline
error. The decisive addition, a quarter-resolution geometry encoding
volume with fail-soft fusion, reached 0.261 pixels, the campaign best,
at 0.041 million added parameters. Step and batch budgets differ
slightly across runs, so the table is an engineering record, not a
controlled comparison; the next subsection's matched studies
re-validate these conclusions.

\begin{table}[!htbp]
\centering
\caption{\textbf{Architecture ablation of the recurrent family.} 20-pair overfit harness at $384\times 640$; steps and batch differ slightly across runs, so the table is an engineering record. The geometry-volume variant became the final model.}
\label{tab:ablation_gev4line}
\small
\begin{tabular}{lccccc}
\toprule
Variant & Params (M) & EPE & Median & bad-1 (\%) & Latency (ms)\\
\midrule
recurrent context (base) & 2.831 & 0.424 & 0.167 & 7.62 & 10.7\\
+ geometry-guided state & 2.919 & 0.459 & 0.218 & 8.16 & 36.7\\
+ update gate & 2.921 & 0.473 & 0.246 & 7.75 & 12.2\\
gate, YOLO26n encoder & 1.410 & 0.387 & 0.182 & 5.98 & 10.9\\
gate, YOLO26s encoder & 2.921 & 0.319 & 0.127 & 5.24 & 11.9\\
+ sharp tail & 2.950 & 0.322 & 0.136 & \textbf{5.21} & 14.6\\
+ high-resolution refine & 3.106 & 0.351 & 0.186 & 5.18 & 16.1\\
+ selective recurrent unit & 3.398 & 0.474 & 0.255 & 8.20 & 14.8\\
+ raw 1/4 initialization & 2.948 & 0.455 & 0.195 & 10.64 & 20.7\\
\textbf{+ GEV and fail-soft fusion} & 2.962 & \textbf{0.261} & \textbf{0.059} & 5.25 & 24.6\\
\bottomrule
\end{tabular}
\end{table}

\FloatBarrier
\subsection{Training-Recipe and Efficiency Studies}
Four further controlled studies (Figure \ref{fig:ablations}) justify
the training recipe and efficiency pass of the final model. The
training recipe, not any architectural block, was the largest accuracy
lever: in a matched nine-arm comparison, the OpenStereo augmentation
triplet lowered the best validation end-point error from 2.778 to
1.921 pixels, a 31 percent reduction well above the noise floor. The
edge-inference optimization pass (Table \ref{tab:ablation_efficiency})
held accuracy while cutting latency by roughly a third:
exact-equivalence fixes plus the narrow geometry volume brought
latency from 46.6 to 30.2 milliseconds, with end-point error unchanged
within noise (2.908 to 2.906 pixels). An exact-equivalence class
(zero-copy views, batched lookups, hoisted static channels, fused
heads) is bitwise or mathematically identical, so its accuracy deltas
are pure noise; the narrow geometry band, the one accuracy-relevant
change, was validated separately as neutral. The boundary-sharpness
study (Table \ref{tab:ablation_blurfix}) selected plane rendering,
which cut the bad-0.5 rate from 55.36 to 42.91 percent, the sharpest
and fastest result of the study, since rendering from planes replaces
a costlier path. Adding the bimodal head or top-k initialization on
top of plane rendering worsened results (44.40 and 46.58 percent
bad-0.5): the fixes anti-synergize, so only plane rendering was
adopted, a negative result that constrains future work. Table
\ref{tab:ablation_native} reports the input-protocol study on both
axes: training on native crops cut the native-axis end-point error
from 6.672 to 2.869 pixels versus the legacy global-resize protocol,
at no cost on the resized axis, so the native-crop protocol was
adopted for the production run.

\begin{table}[!htbp]
\centering
\caption{\textbf{Efficiency ablation.} 100-pair matched harness, 9k steps, identical seed and data; latency at $384\times 640$, batch 1 on the research harness. These latencies index the optimization on the research harness, not the deployment path; the RTX 3050 full-precision measurement of the same pass (106.7 to 61.4 ms, $1.74\times$) is in Table \ref{tab:efficiency_findings}.}
\label{tab:ablation_efficiency}
\small
\begin{tabular}{lccccc}
\toprule
Variant & Params (M) & EPE & bad-1 (\%) & D1-all (\%) & Latency (ms)\\
\midrule
base (gev4) & 2.962 & 2.908 & 30.32 & 16.31 & 46.6\\
+ exact-equivalence fixes & 2.962 & 3.021 & 32.09 & 17.79 & 35.8\\
+ narrow GEV band & 2.962 & 2.906 & 32.29 & 16.88 & 30.2\\
\bottomrule
\end{tabular}
\end{table}

\begin{table}[!htbp]
\centering
\caption{\textbf{Boundary-sharpness fixes and their composition.} 500-pair windowed split; the composition rows show that combining fixes is worse than plane rendering alone.}
\label{tab:ablation_blurfix}
\small
\begin{tabular}{lcccc}
\toprule
Variant & EPE & bad-0.5 (\%) & D1-all (\%) & Latency (ms)\\
\midrule
control (narrow) & 2.006 & 55.36 & 12.57 & 22.2\\
+ bundle1 (top-k init) & 1.982 & 49.95 & 12.14 & 21.9\\
+ bimodal aux head & 1.779 & 43.26 & 11.53 & 25.2\\
\textbf{+ plane rendering (adopted)} & 1.826 & \textbf{42.91} & 12.19 & 15.0\\
\midrule
plane + bimodal & 1.847 & 44.40 & 11.62 & 25.5\\
plane + bimodal + bundle1 & 1.853 & 46.58 & 11.93 & 17.2\\
\bottomrule
\end{tabular}
\end{table}

\begin{table}[!htbp]
\centering
\caption{\textbf{Input-protocol ablation.} Dual-axis final evaluation on the same 50 validation frames; each arm evaluated on both axes. Native crops win the native axis outright at no accuracy cost on the legacy resized axis.}
\label{tab:ablation_native}
\small
\begin{tabular}{lcccccc}
\toprule
& \multicolumn{3}{c}{Native axis ($960\times 540$)} & \multicolumn{3}{c}{Resized axis ($640\times 384$)}\\
\cmidrule(lr){2-4}\cmidrule(lr){5-7}
Arm & EPE & bad-1 & D1-all & EPE & bad-1 & D1-all\\
\midrule
global resize (legacy) & 6.672 & 58.69 & 35.81 & 1.783 & 29.91 & 11.44\\
\textbf{native crop (adopted)} & \textbf{2.869} & \textbf{36.75} & \textbf{16.47} & 1.783 & \textbf{29.30} & \textbf{10.71}\\
native full frame & 2.967 & 38.68 & 16.74 & 2.557 & 36.45 & 16.36\\
\bottomrule
\end{tabular}
\end{table}

\begin{figure}[!htbp]
\centering
\includegraphics[width=\textwidth]{fig_4_7_ablations}
\caption{\textbf{Ablation summary.} (a) OpenStereo augmentation triplet (matched 9-arm comparison, best validation EPE). (b) The edge-inference optimization pass (100-pair harness). (c) Boundary blur fixes and their compositions (500-pair split, bad-0.5). (d) Input protocol on the native evaluation axis.}
\label{fig:ablations}
\end{figure}

\FloatBarrier
\section{Results}

\subsection{In-Domain Accuracy: SceneFlow Test}
The trained model was evaluated on the full FlyingThings3D test set,
4,370 pairs at native 960 by 540 resolution, using the best checkpoint.
Table \ref{tab:sceneflow_results} reports 0.78 pixel end-point error,
3.40 percent D1-all, and a median error of 0.13 pixels; four frames with
no valid ground-truth pixel after the disparity cutoff are excluded from
the macro-average. Table \ref{tab:sceneflow_comparison} adds published
methods for context, though not strictly protocol-identical: each
published EPE follows its own paper's SceneFlow protocol, while this
value comes from the full test set at native resolution. The model falls
between the efficient and iterative families: more accurate than the
classical 3D-cost-volume baseline PSMNet, close to the efficient
LightStereo-S, and a fraction of the parameters of the iterative and
foundation-scale systems. The metric spread in Table
\ref{tab:sceneflow_results} tells its own story: the median of 0.13
pixels is six times smaller than the mean of 0.78, and the RMSE of 3.64
is nearly five times larger, so most pixels are highly accurate while a
small population of large outliers sets the average. The bad-rate ladder
agrees: errors above half a pixel affect 15.3 percent of pixels, but
above three pixels only 4.0 percent, so most failures are moderate and
catastrophic ones rare. Figure \ref{fig:sceneflow_qualitative} shows
where those outliers live, occlusion boundaries and thin structures,
consistent with recurrent reasoning that runs at quarter resolution and
an upsampler that can only recombine, not invent, disparity.

\begin{table}[!htbp]
\centering
\caption{\textbf{Full FlyingThings3D test-set results.} All eight metrics over 4{,}370 pairs at the native $960\times 540$ axis, best checkpoint (step 53k); four frames with no valid ground-truth pixels are excluded from the macro-average.}
\label{tab:sceneflow_results}
\small
\begin{tabular}{cccccccc}
\toprule
EPE (px) & RMSE & Median & bad-0.5 (\%) & bad-1 (\%) & bad-2 (\%) & bad-3 (\%) & D1-all (\%)\\
\midrule
0.781 & 3.64 & 0.130 & 15.32 & 8.92 & 5.34 & 4.00 & 3.40\\
\bottomrule
\end{tabular}
\end{table}

\begin{table}[!htbp]
\centering
\caption{\textbf{Comparison with published methods on SceneFlow.} Published EPE values follow each paper's standard SceneFlow protocol; the StereoLite value is the full FlyingThings3D test set at the native axis and is therefore not strictly protocol-identical. Latency hardware differs across papers.}
\label{tab:sceneflow_comparison}
\small
\begin{tabular}{lcccc}
\toprule
Method & Params (M) & SF EPE (px) & Latency (ms) & Hardware\\
\midrule
PSMNet \cite{chang2018psmnet} & 5.2 & 1.09 & 410 & Titan Xp\\
HITNet (L) \cite{tankovich2021hitnet} & 0.97 & 0.43 & 54 & Titan V\\
CoEx \cite{bangunharcana2021coex} & 2.72 & 0.69 & 27 & RTX 2080Ti\\
RAFT-Stereo \cite{lipson2021raftstereo} & 11.23 & 0.61 & 380 & RTX 6000\\
IGEV-Stereo \cite{xu2023igev} & 12.60 & 0.47 & 180 & RTX 3090\\
LightStereo-S \cite{guo2025lightstereo} & 3.44 & 0.73 & 17 & RTX 3090\\
FoundationStereo \cite{wen2025foundationstereo} & ${\sim}340$ & 0.34 & ${\sim}470$ & RTX 4090\\
\midrule
StereoLite (ours) & 2.96 & 0.78 & 49.8 (fp16) & RTX 3050 laptop\\
\bottomrule
\end{tabular}
\end{table}

\begin{figure}[!htbp]
\centering
\includegraphics[width=\textwidth]{fig_4_3_sceneflow_qualitative}
\caption{\textbf{Qualitative results on held-out FlyingThings3D scenes.} Left image, ground-truth disparity, prediction at the best checkpoint (step 53k), and absolute error clipped at 3 px. Residual error concentrates on occlusion boundaries and sub-pixel structure.}
\label{fig:sceneflow_qualitative}
\end{figure}

\FloatBarrier
\subsection{Zero-Shot Cross-Domain: Middlebury 2014}
Cross-domain accuracy is the real test for an edge model: one accurate
only on its training domain has little field use. The trained model was
evaluated zero-shot on the 23 perfect-set scenes of Middlebury 2014,
unseen during training, under a protocol identical for every method in
Table \ref{tab:mb14_comparison}, reaching 1.71 pixel end-point error and
10.9 percent D1-all, the smallest model in the comparison at 2.96 million
parameters. A gap remains to LiteAnyStereo (6.9 percent) and
sixteen-iteration IGEV-Stereo (5.0 percent), both larger models;
LiteAnyStereo also adds a three-stage knowledge-distillation pipeline the
supervised-only recipe here skips. Figure \ref{fig:mb14_perscene} shows
sixteen of the 23 scenes within the reference band, with the aggregate
dragged up by a few hard scenes; Table \ref{tab:mb14_perscene_full} lists
the per-scene split. The easiest
scenes are ordinary furnished rooms: Adirondack at 2.4 percent D1-all,
Recycle at 2.9, and Storage at 4.2, comparable to the reference models.
The hardest four alone raise the mean by several points: Jadeplant at
34.9 percent, Flowers at 30.3, Mask at 24.7, and Pipes at 17.9. These
scenes share dense thin repeated structures such as leaves and rods,
plus true disparities near or beyond the 192-pixel evaluation cap,
exactly where a quarter-resolution, bounded-search prediction performs
worst. Figure \ref{fig:mb14_qualitative} shows both ends of that
spectrum. The model's cross-domain failure is not diffuse: it is
localized to an identifiable scene type, which the distillation
direction of Chapter 9 targets directly.

\begin{table}[!htbp]
\centering
\caption{\textbf{Zero-shot Middlebury 2014 comparison.} 23 perfect-set scenes, $384\times 640$ inference, identical protocol across all rows (disparity ground truth scaled to the inference width, values above 192 px masked). Reference rows from the repository's matched-protocol evaluations of the public checkpoints \cite{jing2025liteanystereo,xu2023igev}. Latency is the per-scene median on a cloud T4 in fp32 eager mode.}
\label{tab:mb14_comparison}
\small
\resizebox{\textwidth}{!}{%
\begin{tabular}{lcccccccc}
\toprule
Model & Params (M) & T4 (ms) & EPE & bad-0.5 & bad-1 & bad-2 & bad-3 & D1-all (\%)\\
\midrule
\textbf{StereoLite (ours)} & 2.96 & 86 & \textbf{1.71} & 39.6 & 23.9 & 14.5 & 11.2 & \textbf{10.9}\\
LiteAnyStereo \cite{jing2025liteanystereo} & 7.60 & 64 & 1.17 & 35.1 & 19.7 & 10.5 & 7.4 & 6.9\\
IGEV-Stereo (16 it.) \cite{xu2023igev} & 12.60 & 305 & 0.86 & 22.2 & 12.4 & 7.2 & 5.2 & 5.0\\
\bottomrule
\end{tabular}}
\end{table}

\begin{figure}[H]
\centering
\includegraphics[width=\textwidth]{fig_4_4_mb14_perscene}
\caption{\textbf{Per-scene zero-shot D1-all on Middlebury 2014.} 23 perfect-set scenes sorted by error; horizontal lines mark the aggregate means of our model and the two reference models under the identical protocol.}
\label{fig:mb14_perscene}
\end{figure}

\begin{figure}[H]
\centering
\includegraphics[width=0.9\textwidth]{fig_4_5_mb14_qualitative}
\caption{\textbf{Zero-shot qualitative results on Middlebury 2014.} The three easiest and two hardest scenes by D1-all: left image, ground truth, and our prediction (turbo colormap, per-scene scale, invalid ground truth dark). The hardest scenes combine thin structures with disparities near the evaluation cap.}
\label{fig:mb14_qualitative}
\end{figure}

\FloatBarrier
\subsection{Robustness to Rectification Imperfection}
\label{sec:rectification}
A second thesis objective is tolerance to imperfect rectification. The
trained model was evaluated on a 400-pair test subset with the right
image shifted vertically by a controlled offset and ground truth
unchanged, emulating residual vertical misalignment between views. Table
\ref{tab:rectification} and Figure \ref{fig:rectification} report the
result. The zero-offset baseline of 1.03 pixels is measured on this
400-pair perturbation subset rather than the full 4,370-pair test set, so
it sits above the 0.78-pixel headline of Table
\ref{tab:sceneflow_results}; the quantity of interest here is the rise
under misalignment, not the absolute level. Accuracy degrades gracefully
across the practical range. At a half-pixel offset, end-point error rises
only from 1.03 to 1.22 pixels
and D1-all from 4.3 to 4.6 percent; at a full pixel, error stays at 1.53
pixels with 6.2 percent D1-all. A calibrated rig typically holds residual
vertical error below one pixel, so accuracy survives this range; the
tolerance comes from the refinement stage's local correlation lookup,
which checks a small neighborhood around each match instead of assuming
a horizontal epipolar line. Beyond this range degradation is steep: a
two-pixel offset raises D1-all to 15.8 percent, a four-pixel offset to
41.2 percent, since the network lacks explicit vertical search and large
misalignment pushes matches outside the lookup window. This is measured
tolerance to small misalignment, not immunity to arbitrary
de-rectification.

\begin{table}[!htbp]
\centering
\caption{\textbf{Accuracy under vertical rectification error.} The right image is shifted vertically by the stated offset on a 400-pair FlyingThings3D test subset at the native axis; the disparity ground truth is unchanged. Accuracy is retained across the sub-pixel to one-pixel range that a calibrated rig produces.}
\label{tab:rectification}
\small
\begin{tabular}{lcccc}
\toprule
Vertical offset (px) & EPE (px) & bad-1 (\%) & bad-2 (\%) & D1-all (\%)\\
\midrule
0.0 & 1.03 & 10.28 & 6.63 & 4.29\\
0.5 & 1.22 & 14.62 & 7.58 & 4.60\\
1.0 & 1.53 & 26.27 & 11.84 & 6.19\\
2.0 & 2.43 & 47.73 & 27.71 & 15.83\\
4.0 & 5.07 & 71.59 & 55.68 & 41.17\\
\bottomrule
\end{tabular}
\end{table}

\begin{figure}[!htbp]
\centering
\includegraphics[width=0.72\textwidth]{fig_4_9_rectification}
\caption{\textbf{Accuracy under vertical rectification error.} End-point error and D1-all against the right-image vertical offset; the shaded band marks the sub-pixel to one-pixel residual a calibrated rig produces, within which accuracy is largely retained.}
\label{fig:rectification}
\end{figure}

The synthetic sweep is controlled with known ground truth, but the
condition it emulates occurs for real on a low-cost rig. To confirm this
on physical data, the model was run on 997 rectified pairs from the
project's stereo rig, a commodity camera whose factory rectification
leaves residual vertical misalignment, and compared against the
FoundationStereo teacher's pseudo-ground-truth on the same pairs.
Agreement is close: mean end-point error is 1.45 pixels and D1-all is
11.1 percent, matching the in-domain synthetic error despite the real
rig's imperfect rectification. Since both see the same misaligned input,
this figure measures agreement with a strong reference rather than
absolute error, corroborating rather than replacing the controlled sweep.
Together, the two results confirm that the model withstands the
rectification imperfections of an inexpensive camera.

\FloatBarrier
\subsection{Efficiency and Deployment Measurements}
Table \ref{tab:latency} reports the measured efficiency. On the RTX 3050
laptop GPU at batch one and 384 by 640 resolution, the optimized model ran
in 49.8 milliseconds at half precision and 61.4 milliseconds at full
precision, with inference working sets of 0.26 and 0.35 gigabytes. On the
target device, a Jetson Orin Nano, the optimized eight-bit model measured
36.3 milliseconds per frame on-device, or 27.5 frames per second at the
same resolution. This nears the 30 frames per second real-time target,
confirming interactive rates on the intended edge accelerator, not just
the laptop GPU. The on-device working set stays well within the
platform's memory budget.

\begin{table}[!htbp]
\centering
\caption{\textbf{Latency and memory.} RTX 3050 Laptop values are batch 1, $384\times 640$, eager PyTorch. The Jetson Orin Nano row is a direct on-device measurement of the optimized eight-bit TensorRT model; its memory figure is the estimated on-device working set.}
\label{tab:latency}
\small
\begin{tabular}{lcccc}
\toprule
Configuration & Latency (ms) & FPS & Inference memory (GB)\\
\midrule
RTX 3050, fp32 & 61.4 & 16.3 & 0.35\\
RTX 3050, fp16 & 49.8 & 20.1 & 0.26\\
Jetson Orin Nano, INT8 TensorRT & 36.3 & 27.5 & ${\sim}0.15$\\
\bottomrule
\end{tabular}
\end{table}

\FloatBarrier
\subsection{Qualitative Results}
Visual inspection supports quantitative results. On the synthetic test
set, Figure \ref{fig:sceneflow_qualitative} shows accurate reconstruction
of broad surfaces, with error confined to occlusion boundaries. On the
unseen Middlebury scenes, Figure \ref{fig:mb14_qualitative} shows clean
reconstruction of easy scenes, with failure localized to thin repeated
structures, not spread across the image. Chapter 3's plane-rendering step
shows at the boundaries: predicted disparity follows object edges more
closely than bilinear upsampling. Figure \ref{fig:camera} extends this to
real data: four indoor pairs from the calibrated stereo rig, none used in
training or fine-tuning, yield plausible disparity maps that cleanly
separate foreground from background, and the model transfers to a real
camera without adaptation.

\begin{figure}[!htbp]
\centering
\includegraphics[width=\textwidth]{fig_4_6_camera}
\caption{\textbf{Real-camera qualitative results.} Four indoor pairs from the project's stereo rig ($1280\times 720$, rectified, no fine-tuning): left image and the model's predicted disparity at $384\times 640$. No ground truth exists for these pairs; the panels demonstrate zero-shot visual plausibility on real data.}
\label{fig:camera}
\end{figure}

\FloatBarrier
\subsection{3D Reconstruction of Real Scenes}
This section presents the promised reconstruction. Disparity from three
real indoor scenes was back-projected to metric point clouds via the
calibrated relations of \eqref{eq:backprojection}, cleaned with a
statistical outlier filter and voxel downsampling per Section
\ref{sec:reconstruction_method}, and rendered from three viewpoints
(Figure \ref{fig:reconstruction}). The reconstructions are metrically
plausible. In the first scene, two fire extinguishers stand proud of the
wall and door in the oblique and top-down views. In the second, the
corridor recovers a smooth depth gradient to the bright far end. In the
third, the metal cabinet sits in front of the back wall at its correct
standoff. The flying points along depth boundaries are the
three-dimensional footprint of the same occlusion-edge error in the
error maps of Figure \ref{fig:sceneflow_qualitative}; the depth
quantization on distant surfaces reflects depth's inverse relation to
disparity: a fixed sub-pixel disparity step spans a growing depth
interval with distance. These pairs come from the same low-cost rig
whose imperfect rectification is examined in Section
\ref{sec:rectification}; no fine-tuning was done on the rig's data, so
the clouds are the zero-shot output of the checkpoint evaluated
throughout this chapter.

\begin{figure}[H]
\centering
\includegraphics[width=\textwidth]{fig_4_10_reconstruction}
\caption{\textbf{3D reconstruction of real scenes from predicted disparity.} Three indoor pairs from the project's low-cost stereo rig: the left image and the colored point cloud rendered from the camera viewpoint, an oblique viewpoint, and top-down. Clouds are back-projected through the calibrated rig geometry, outlier-filtered, and voxel-downsampled; no fine-tuning on the rig's data was performed.}
\label{fig:reconstruction}
\end{figure}

\FloatBarrier
\section{Discussion}
The results answer the four research questions from Chapter 3. The
model reaches 0.78 pixel end-point error within a few million
parameters, answering the first. It generalizes zero-shot to an unseen
domain at 10.9 percent D1-all, short of the distillation-trained
references, answering the second. It runs interactively on a laptop GPU
and hit 27.5 frames per second on the Jetson Orin Nano, near the
real-time target, answering the third. The fourth, tolerance to imperfect
rectification, is addressed by the perturbation study in Section
\ref{sec:rectification}. Figure \ref{fig:pareto_ours} places the model on
the zero-shot accuracy-efficiency front, at the low-parameter end below
the larger distillation-trained and iterative references. The comparison
bears on an open question about cross-domain generalization: LiteAnyStereo,
the strongest efficient reference, is itself non-iterative and owes its
accuracy to a three-stage distillation recipe rather than to iteration,
which suggests the training pipeline carries more of the cross-domain gain
than architectural iteration does, though this study does not isolate the
two factors. Limitations: results come from
a single run, so run-to-run variance is unknown; the model was not
submitted to the KITTI leaderboard; and it lacks left-right consistency
or occlusion reasoning. These limits frame the future directions in
Chapter 9.

\begin{figure}[!htbp]
\centering
\includegraphics[width=0.78\textwidth]{fig_4_8_pareto_ours}
\caption{\textbf{StereoLite on the zero-shot accuracy-efficiency front.} Middlebury 2014 D1-all against trainable parameters; all three points share the identical 23-scene protocol of Table \ref{tab:mb14_comparison}.}
\label{fig:pareto_ours}
\end{figure}

\FloatBarrier
\section{Summary}
This chapter evaluates the model on every objective axis. In domain, it
reached 0.78 pixel end-point error and 3.40 percent D1-all on the
FlyingThings3D test set. Zero-shot on unseen Middlebury 2014 scenes, it
reached 10.9 percent D1-all under a protocol identical to the reference
models, as the smallest network in that comparison; remaining failures
cluster in scenes with thin repeated structures and disparities near the
cap. Under controlled vertical rectification error, end-point error rose
from 1.03 to 1.53 pixels across the sub-pixel to one-pixel range a
calibrated rig produces; the real-rig agreement study confirmed that
tolerance. On the Jetson Orin Nano, the optimized eight-bit model hit
36.3 milliseconds per frame, near the 30 frames per second real-time
goal; real-camera disparity maps and metric point-cloud reconstructions
confirmed the checkpoint transfers to physical data without adaptation.
Chapter 9 draws together the conclusions, limitations, and future
directions these results support.
